% BHU PYQ -- Manually transcribed & error-corrected
% Paper: GRB-403A: Basics of Remote Sensing
% Exam: B.Sc. (Hons.) Semester IV Examination, 2023-24

\documentclass[11pt,a4paper]{article}
\usepackage[a4paper,top=2.5cm,bottom=2.5cm,left=2.8cm,right=2.8cm,headheight=14pt]{geometry}
\usepackage{amsmath,amssymb}
\usepackage[T1]{fontenc}
\usepackage[utf8]{inputenc}
\usepackage{lmodern,microtype,fancyhdr,enumitem,hyperref,lastpage,parskip}

\hypersetup{colorlinks=true,urlcolor=black,linkcolor=black,
  pdftitle={GRB-403A Basics of Remote Sensing -- BHU B.Sc. Sem IV 2023-24}}
\pagestyle{fancy}\fancyhf{}
\renewcommand{\headrulewidth}{0.4pt}\renewcommand{\footrulewidth}{0.4pt}
\fancyhead[L]{\small   \textit{Banaras Hindu University}}
\fancyhead[R]{\small   \textit{GRB-403A: Basics of Remote Sensing}}
\fancyfoot[C]{\small Page  \thepage\ of \pageref{LastPage}}


\newlist{parts}{enumerate}{2}
\setlist[parts,1]{label=(\alph*),leftmargin=*,itemsep=5pt,topsep=3pt}

\newcommand{\pts}[1]{\hfill   \textit{[#1]}}


\begin{document}

\begin{center}
  {\large\bfseries BANARAS HINDU UNIVERSITY}\\[2pt]
  {
\normalsize B.Sc. (Hons.) --- Semester IV Examination, 2023-24}\\[6pt]
  \rule{0.8\linewidth}{0.5pt}\\[6pt]
  {\large\bfseries Paper No. GRB-403A: Basics of Remote Sensing}\\[4pt]
  {
\normalsize Time: 3 Hours \hfill Maximum Marks: 70}\\[2pt]
  {\small Ref. No.: 074555}
\end{center}
\rule{\linewidth}{0.4pt}
\smallskip



\noindent   \textit{Instructions: Answer   \textbf{five} questions in all including Question 1 which is compulsory. Remaining questions are of equal marks.}
\bigskip




\noindent   \textbf{Question 1.} Write short notes on the following: \pts{20}

\begin{parts}
  \item Rayleigh scattering
  \item Characteristics of oblique aerial photographs.
  \item Electromagnetic radiation.
  \item Sun-synchronous satellites.
  \item Use of association in visual image interpretation.
\end{parts}

\medskip


\noindent   \textbf{Question 2.} Explain the interaction mechanism of electromagnetic radiations with the atmosphere. \pts{12.5}

\medskip


\noindent   \textbf{Question 3.} Explain the characteristics of different platforms used in remote sensing. \pts{12.5}

\medskip


\noindent   \textbf{Question 4.} Write a note on digital image processing techniques. \pts{12.5}

\medskip


\noindent   \textbf{Question 5.} Describe the steps and procedures involved in image interpretation. \pts{12.5}

\medskip


\noindent   \textbf{Question 6.} Identify and discuss the various applications of remote sensing in management of forest resources. \pts{12.5}

\medskip


\noindent   \textbf{Question 7.} Discuss the characteristics of any three sensors used in Indian remote sensing satellites. \pts{12.5}

\medskip


\noindent   \textbf{Question 8.} Explain the interaction mechanism of electromagnetic radiations with Earth's surface features. \pts{12.5}

\vfill\rule{\linewidth}{0.4pt}

\begin{center}\small
\href{https://bhu-exam.vercel.app/}{Click here to visit our website}\\[4pt]
\href{https://me-aryan.vercel.app/}{Click here to visit our contributor}\\[4pt]
\href{https://quashan-me.vercel.app/}{Click here to visit our contributor}
\end{center}
\end{document}
